\chapter[2014 --- O'Keefe et al.]{2014: John O'Keefe, May-Britt Moser, Edvard I. Moser}
\label{ch:2014_OKeefe}

\section*{Main Contribution}

\textit{Discovered grid cells and place cells that constitute a positioning system in the brain, explaining how the brain creates a cognitive map of the environment.}

\section{Official Citation}

for their discoveries of cells that constitute a positioning system in the brain

\section{Background and Historical Context}

The 2014 Nobel Prize in Physiology or Medicine recognized John O'Keefe, May-Britt Moser, and Edvard I. Moser ``for their discoveries of cells that constitute a positioning system in the brain.'' Their work revealed how the brain creates a map of the surrounding environment and enables navigation, solving one of the most intriguing mysteries in neuroscience: how do we know where we are and how do we find our way from one place to another?

The ability to navigate through space is fundamental to the survival of virtually all mobile organisms, from insects to humans. Animals must be able to recognize their environment, remember locations of food, water, shelter, and mates, and plan efficient routes between different locations. The neural mechanisms underlying spatial navigation had been a subject of intense interest in neuroscience since the pioneering work of Edward Tolman in the 1940s, who proposed the concept of ``cognitive maps''---internal representations of the spatial layout of the environment that animals use to navigate.

In the 1970s, the British psychologist John O'Keefe, working at University College London, made a pioneering discovery that provided the first neural evidence for cognitive maps. O'Keefe was studying the electrical activity of neurons in the hippocampus, a brain structure that had been implicated in memory and spatial processing. Using implanted electrodes to record the activity of individual neurons in the hippocampi of rats as they explored enclosed environments, O'Keefe observed that certain neurons fired selectively when the rat occupied specific locations in the environment. He named these neurons ``place cells'' and proposed that their activity represented a neural map of the environment---the cognitive map that Tolman had hypothesized decades earlier.

The discovery of place cells was initially controversial. Some critics argued that the firing of place cells might reflect non-spatial factors, such as sensory stimuli, motor activity, or motivational states, rather than the animal's location per se. However, O'Keefe and his colleagues conducted extensive experiments demonstrating that place cell activity was specifically tied to the animal's spatial location and was not simply a reflection of other variables. They showed that place cells had stable firing fields that persisted over days and weeks, that different place cells encoded different locations, and that the ensemble activity of many place cells could accurately predict the animal's position in the environment.

\section{The Discovery/Contribution}

While O'Keefe's discovery of place cells in the hippocampus was a landmark achievement, it left a fundamental question unanswered: how are place cell representations formed and updated as the animal moves through the environment? The place cells of the hippocampus appeared to provide a readout of the animal's position, but the input signals that drove place cell activity---the sensory information that was used to construct the spatial map---remained unknown.

This question was addressed by May-Britt Moser and Edvard I. Moser, a married couple from Norway who were working at the Norwegian University of Science and Technology in Trondheim. The Mosers began their investigations in the early 2000s, building on the foundation laid by O'Keefe's work and using newly developed recording techniques that allowed them to simultaneously monitor the activity of large numbers of neurons in freely behaving animals.

In a landmark study published in 2005, the Mosers discovered a new type of spatially tuned neuron in the medial entorhinal cortex (MEC), a brain region that provides input to the hippocampus. These neurons, which they named ``grid cells,'' fired at multiple locations in the environment, arranged in a strikingly regular hexagonal pattern. As the animal moved through the environment, grid cells fired at regularly spaced positions, creating a periodic grid-like representation of space that covered the entire environment.

The discovery of grid cells was a revelation. The hexagonal firing pattern of grid cells was unlike anything previously observed in the brain and suggested that the entorhinal cortex contained a metric representation of space---a neural coordinate system that could be used to calculate distances and directions. Grid cells appeared to provide the spatial information that was used by place cells in the hippocampus to construct detailed maps of the environment.

The Mosers subsequently identified additional cell types in the entorhinal cortex and surrounding areas that contributed to the brain's navigation system. Head direction cells, which fired when the animal faced a specific direction regardless of its location, were found in the presubiculum and provided information about the animal's heading. Border cells (also called boundary cells), which fired when the animal was near the walls or boundaries of the environment, were found in the entorhinal cortex and provided information about the animal's position relative to environmental boundaries. Speed cells, which fired at a rate proportional to the animal's running speed, were also identified in the entorhinal cortex.

The identification of these diverse cell types revealed that the brain's navigation system was far more complex and sophisticated than initially appreciated. Rather than a simple map, the navigation system appeared to be a modular computational network in which different cell types provided different types of spatial information that were integrated to produce a coherent representation of space. This modular organization was reminiscent of the computational architectures used in robotic navigation systems, where different sensors provide complementary information that is combined to produce a robust estimate of position and orientation.

The Mosers' technical innovations were also crucial to their discoveries. They developed large-scale recording techniques that allowed them to simultaneously monitor the activity of hundreds of neurons from multiple brain regions in freely behaving animals. These recordings, which used silicon probes with multiple recording sites, provided notable resolution of the spatial organization of the navigation system and enabled the identification of cell types that would have been difficult to detect using single-unit recording methods.

Together, these cell types---place cells in the hippocampus, grid cells, head direction cells, border cells, and speed cells in the entorhinal cortex and surrounding areas---constituted a comprehensive ``positioning system'' in the brain. This system, sometimes referred to as the ``GPS of the brain'' or the ``hippocampal-entorhinal navigation system,'' provided the neural basis for spatial navigation and the formation of cognitive maps.

The Mosers' work also revealed that the navigation system was organized in a hierarchical manner. Grid cells in different parts of the entorhinal cortex had different grid spacings, with smaller grids in the dorsolateral MEC and larger grids in the ventromedial MEC. This gradient of grid spacings suggested that the navigation system could represent space at multiple scales, from small local environments to large, complex landscapes.

\section{Impact on Modern Medicine}

The discoveries of O'Keefe and the Mosers have had a significant impact on our understanding of brain function and have opened up new avenues for the treatment of neurological and psychiatric disorders.

One of a major clinical implications of the navigation system discoveries relates to Alzheimer's disease. The hippocampus and entorhinal cortex are among the earliest brain regions affected in Alzheimer's disease, and spatial disorientation---difficulty finding one's way in familiar environments---is one of the earliest and most characteristic symptoms of the disease. The understanding of the neural mechanisms of spatial navigation has provided a framework for understanding why spatial disorientation occurs so early in Alzheimer's disease and has suggested potential strategies for early diagnosis and intervention.

Research has shown that the entorhinal cortex is selectively vulnerable to the neurodegenerative process of Alzheimer's disease, with grid cells being among the first neurons to be affected. This vulnerability may explain why spatial navigation deficits are among the earliest clinical manifestations of the disease, often appearing years before memory loss becomes apparent. The development of sensitive tests of spatial navigation performance has been explored as a potential tool for early detection of Alzheimer's disease, potentially enabling earlier diagnosis and intervention.

The navigation system has also been implicated in other neurological and psychiatric conditions. Deficits in spatial navigation have been observed in individuals with epilepsy, schizophrenia, post-traumatic stress disorder, and autism spectrum disorder. Understanding the neural basis of spatial navigation has provided insights into the mechanisms underlying these conditions and has suggested potential therapeutic targets.

In the field of robotics and artificial intelligence, the discovery of grid cells has inspired the development of new algorithms for spatial navigation and mapping. The hexagonal grid pattern of grid cells has been used as the basis for bio-inspired navigation systems that can be applied in autonomous vehicles, drones, and robots. The mathematical properties of grid cell coding have been shown to provide an efficient and robust representation of space that outperforms many artificial navigation algorithms.

The work of O'Keefe and the Mosers has also contributed to our understanding of the relationship between spatial navigation and memory. The hippocampus, which contains place cells, is also critically involved in the formation of episodic memories---memories of specific events and experiences that are tied to a particular time and place. The overlap between the navigation system and the memory system suggests that spatial information may serve as a framework for organizing and storing memories, a hypothesis that is supported by extensive experimental evidence.

John O'Keefe was awarded half of the 2014 Nobel Prize in Physiology or Medicine, while May-Britt Moser and Edvard I. Moser shared the other half, ``for their discoveries of cells that constitute a positioning system in the brain.'' Their work has fundamentally changed our understanding of how the brain processes spatial information and has had a lasting impact on neuroscience and medicine.

\section{Legacy}

The legacy of O'Keefe, May-Britt Moser, and Edvard I. Moser's work extends across multiple areas of neuroscience and medicine. Their discoveries have revealed a fundamental neural system that enables organisms to navigate through their environment, and this system has proven to be far more complex and sophisticated than initially appreciated.

John O'Keefe's original discovery of place cells in the 1970s established the hippocampus as a key brain structure for spatial processing and provided the first neural evidence for cognitive maps. His work inspired decades of research on the neural mechanisms of spatial navigation and memory, and his place cell theory continues to be a central framework in hippocampal neuroscience.

May-Britt and Edvard Moser's discovery of grid cells and additional cell types in the entorhinal cortex has revealed the computational principles underlying the brain's navigation system. The grid cell, with its elegant hexagonal firing pattern, has become one of the most iconic and widely studied neural representations in neuroscience. The Mosers' work has also established the entorhinal cortex as a major computational hub that processes and integrates multiple types of spatial information before transmitting it to the hippocampus.

The navigation system discoveries have also contributed to a broader understanding of how the brain represents and processes information. The hierarchical organization of the navigation system, with different cell types providing different types of spatial information at different scales, provides a model for understanding how the brain constructs complex representations from simple computational elements. This principle has been applied to the study of other sensory and cognitive systems, contributing to a more general understanding of neural coding and computation.

Today, research on the brain's navigation system continues to be one of the most active and productive areas of neuroscience. New techniques, including optogenetics, calcium imaging, and large-scale neural recording, are being used to study the navigation system with notable precision and resolution. The legacy of O'Keefe and the Mosers' work provides the foundation for these ongoing investigations and continues to inspire new discoveries about how the brain processes spatial information.


\section{Modern Developments}

O'Keefe, and the Mosers' navigation work has led to modern spatial neuroscience. Understanding of grid cells has informed development of GPS-like systems for robotics. Understanding of spatial memory has led to research on Alzheimer's disease, where spatial disorientation is an early symptom. Understanding of hippocampal circuits has informed development of deep brain stimulation for epilepsy. Understanding of place cells has inspired AI navigation algorithms.
